\documentclass[a4,11pt]{aleph-notas}

% -- Paquetes adicionales
\usepackage{aleph-comandos}
\usepackage{booktabs}
\usepackage{pdflscape}

% -- Datos   
\institucion{Facultad de Ciencias Exactas, Naturales y Ambientales}
\carrera{Catálogo STEM}
\asignatura{Álgebra lineal}
\tema{Video exposición no. 1: El Teorema de Rouché-Fröbenius}
\autor{Andrés Merino}
\fecha{Periodo 2026-1}

\logouno[0.16\textwidth]{Logos/logoPUCE_04_ac}
\definecolor{colortext}{HTML}{1957d1}
\definecolor{colordef}{HTML}{1957d1}
\fuente{montserrat}


% -- Otros comandos




\begin{document}

\encabezado


%%%%%%%%%%%%%%%%%%%%%%%%%%%%%%%%%%%%%%%%
\section{Resultado de aprendizaje}
%%%%%%%%%%%%%%%%%%%%%%%%%%%%%%%%%%%%%%%%

\begin{itemize}
\item 
    Esta actividad evalúa si el estudiante ({Criterio 1.2:}) identifica los elementos de los sistemas de ecuaciones lineales y la clasificación de estos por su tipo de soluciones.
\end{itemize}

%%%%%%%%%%%%%%%%%%%%%%%%%%%%%%%%%%%%%%%%
\section{Indicaciones}
%%%%%%%%%%%%%%%%%%%%%%%%%%%%%%%%%%%%%%%%

\begin{itemize}
\item 
    El estudiante deberá elaborar diapositivas profesionales y grabar un video presentándolas, con el rostro visible durante toda la exposición.
\item 
    No se aceptará un video en el que el estudiante únicamente lea las diapositivas; se deberá evidenciar dominio del tema mediante una explicación clara y personal.
\item 
    No se aceptarán presentaciones ni libretos claramente generado por inteligencia artificial y descontextualizado de lo visto en clase.
\item 
    La duración máxima del video será de 7 minutos.
\end{itemize}


%%%%%%%%%%%%%%%%%%%%%%%%%%%%%%%%%%%%%%%%
\section{Descripción}
%%%%%%%%%%%%%%%%%%%%%%%%%%%%%%%%%%%%%%%%

\begin{enumerate}
    \item \textbf{Elaboración de diapositivas:}
    \begin{itemize}
        \item Crear una presentación clara y ordenada sobre los sistemas de ecuaciones lineales y sus tipos de solución.
        \item Explicar los tres posibles casos: solución única, infinitas soluciones y ninguna solución.
        \item Incluir ejemplos concretos y una representación gráfica en tres dimensiones para cada caso.
        \item Presentar el Teorema de Rouché-Frobenius y su interpretación para clasificar las soluciones.
        \item Incluir al menos un ejemplo numérico en el que se aplique el teorema para determinar el tipo de solución.
    \end{itemize}

    \item \textbf{Grabación del video:}
    \begin{itemize}
        \item El estudiante deberá presentarse en cámara durante toda la exposición.
        \item El rostro debe ser visible en todo momento.
        \item No se permite únicamente voz en off ni una grabación en la que solo se lean las diapositivas.
    \end{itemize}

    \item \textbf{Estructura mínima:}
    \begin{itemize}
        \item \textbf{Introducción:} explicación breve sobre qué son los sistemas de ecuaciones lineales y por qué es importante estudiar sus soluciones.
        \item \textbf{Desarrollo:} explicación de los tipos de solución, el teorema y los ejemplos prácticos; junto con la explicación del teorema de Rouché-Frobenius y su aplicación.
        \item \textbf{Cierre:} síntesis final sobre la utilidad del Teorema de Rouché-Frobenius.
    \end{itemize}
\end{enumerate}

%%%%%%%%%%%%%%%%%%%%%%%%%%%%%%%%%%%%%%%%
\section{Entrega}
%%%%%%%%%%%%%%%%%%%%%%%%%%%%%%%%%%%%%%%%

Se deberá entregar en el aula virtual el enlace de \textbf{OneDrive} del video con permisos de acceso habilitados. Si el enlace no permite visualizar el material, la actividad no podrá ser calificada y se asignará \textbf{0}.

%%%%%%%%%%%%%%%%%%%%%%%%%%%%%%%%%%%%%%%%
\section{Rúbrica de calificación}
%%%%%%%%%%%%%%%%%%%%%%%%%%%%%%%%%%%%%%%%

En caso de presentar un video con voz en off, rostro no visible o fuera del rango de tiempo establecido, se reducirá un nivel a cada punto de la rúbrica.

\begin{landscape}
\begin{center}\footnotesize
\begin{tabular}{p{5cm}p{4cm}p{4cm}p{4cm}p{4cm}}
\toprule
\textbf{Indicador} & \textbf{Excelente} & \textbf{Bueno} & \textbf{Aceptable} & \textbf{Insuficiente} \\ \midrule
\textbf{Tipos de soluciones (10 pts)} & Explica con claridad total los tres tipos de solución y aporta ejemplos correctos (10) & Explica los tres tipos con omisiones menores (8) & Describe parcialmente los tipos de solución (6) & Explicación incompleta o incorrecta (0) \\ \midrule
\textbf{Representación gráfica (10 pts)} & Incluye gráficos claros y pertinentes para los casos (10) & Presenta gráficos adecuados con ligeras debilidades (8) & Los gráficos son simples o poco precisos (6) & No incluye gráficos relevantes (0) \\ \midrule
\textbf{Teorema de Rouché-Frobenius (10 pts)} & Explica el teorema, sus condiciones y su uso con precisión (10) & Explicación correcta con detalles faltantes (8) & Menciona el teorema sin desarrollo suficiente (6) & Explicación superficial o errónea (0) \\ \midrule
\textbf{Ejemplos prácticos (10 pts)} & Presenta ejemplos numéricos bien desarrollados y bien interpretados (10) & Incluye ejemplos adecuados con explicación parcial (8) & Los ejemplos son limitados o poco desarrollados (6) & No presenta ejemplos o son incorrectos (0) \\ \midrule
\textbf{Presentación en video (5 pts)}
& Exposición clara y fluida; rostro visible durante toda la presentación (5)
& Exposición clara con ligeras interrupciones; rostro visible (4)
& Exposición poco fluida o lectura excesiva; rostro visible parcial (3)
& Voz en off o rostro no visible / exposición ininteligible (0) \\ \midrule
\textbf{Gestión del tiempo (5 pts)} & 7:00 $\pm$ 30 s (5) & 7:00 $\pm$ 60 s (4) & 7:00 $\pm$ 90 s (3) & Fuera de $\pm$ 90 s (0) \\ \bottomrule
\end{tabular}
\end{center}
\end{landscape}



\end{document}